% ====================================
% CHƯƠNG III - PHẦN 2: UML NGHIỆP VỤ
% Người phụ trách: Người 4
% ====================================






%\FloatBarrier
%\subsection{Đặc tả Use Case}
%\label{subsec:dac_ta_use_case}

% Sử dụng danh sách liệt kê để tránh làm sâu mục lục
%Các Use Case cần đặc tả chi tiết:
%\begin{enumerate}
%    \item \textbf{Đăng nhập hệ thống}
%   \item \textbf{Quản lý Series}
%   \item \textbf{Quản lý Chapter}
%    \item \textbf{Quản lý Task}
%    \item \textbf{Nộp Submission}
%    \item \textbf{Review Chapter}
%    \item \textbf{Xuất bản Manga}
%   \item \textbf{Quản lý Ranking}
%    \item \textbf{Quản lý Notification}
%\end{enumerate}


\section{Thiết kế quy trình nghiệp vụ}
\label{sec:thiet_ke_quy_trinh_nghiep_vu}

\subsection{Swimlane Diagram tổng quát quy trình sáng tác và xuất bản Manga}
\label{subsec:swimlane_manga}
Swimlane Diagram được sử dụng nhằm mô tả tổng quan quy trình sáng tác và xuất bản Manga, đồng thời thể hiện sự tương tác giữa các đối tượng tham gia trong hệ thống gồm Mangaka, Assistant, Tantou Editor và Editorial Board.\\
\indent Quy trình bắt đầu từ giai đoạn đề xuất và xét duyệt Series, tiếp tục với quá trình sáng tác, phân công công việc, kiểm duyệt nội dung và kết thúc bằng việc đánh giá kết quả phát hành để đưa ra quyết định tiếp tục hoặc ngừng xuất bản Series. 
\begin{figure}[H]
    \centering
    \includegraphics[height=0.75\textheight, width=\textwidth, keepaspectratio]{Chapter3_IMAGE/Swimlane_tong_quat_quy_trinh_sang_tac_maga.png}
    \caption{Swimlane Diagram tổng quát quy trình sáng tác và xuất bản Manga}
    \label{fig:swimlane_manga}
\end{figure}

\FloatBarrier
\subsection{Activity Diagram quy trình giao và thực hiện Task}
\label{subsec:activity_task}
Activity Diagram này mô tả quy trình phân công và thực hiện công việc giữa Mangaka và Assistant.\\
\indent Quy trình bắt đầu khi Mangaka tạo Task và giao cho Assistant. Sau khi nhận nhiệm vụ, Assistant thực hiện công việc, gửi Submission lên hệ thống và chờ kết quả kiểm duyệt. Trong trường hợp Submission chưa đạt yêu cầu, Assistant sẽ tiếp tục chỉnh sửa và gửi lại cho đến khi được chấp thuận.
\begin{figure}[H]
    \centering
    \includegraphics[height=0.75\textheight, width=0.8\textwidth, keepaspectratio]{Chapter3_IMAGE/Activity_Diagram_quy_trình_giao_và_thực_hiện_Task.png}
    \caption{Activity Diagram quy trình giao và thực hiện Task}
    \label{fig:activity_task}
\end{figure}

\FloatBarrier
\subsection{Activity Diagram quy trình Review Chapter}
\label{subsec:activity_review}
Activity Diagram này mô tả quy trình kiểm duyệt nội dung Manga giữa Mangaka và Tantou Editor.\\
\indent Sau khi Chapter được hoàn thiện, Tantou Editor tiến hành xem xét nội dung và đưa ra phản hồi. Nếu Chapter chưa đạt yêu cầu, Mangaka sẽ thực hiện chỉnh sửa theo nhận xét. Quy trình được lặp lại cho đến khi Chapter được phê duyệt.
\begin{figure}[H]
    \centering
    \includegraphics[height=0.75\textheight, width=0.8\textwidth, keepaspectratio]{Chapter3_IMAGE/Activity_Diagram_quy trình_Review_Chapter.png}
    \caption{Activity Diagram quy trình Review Chapter}
    \label{fig:activity_review}
\end{figure}

\FloatBarrier
\subsection{Activity Diagram quy trình xuất bản Manga}
\label{subsec:activity_xuat_ban}
Activity Diagram này mô tả quá trình đánh giá và ra quyết định xuất bản đối với một Series Manga.\\
\indent Editorial Board thực hiện cập nhật dữ liệu bình chọn, theo dõi thứ hạng Series và đánh giá hiệu quả phát hành. Dựa trên kết quả đánh giá, hội đồng đưa ra quyết định tiếp tục xuất bản, điều chỉnh kế hoạch phát hành hoặc ngừng phát hành Series.
\begin{figure}[htbp]
    \centering
    \includegraphics[height=0.75\textheight, width=0.8\textwidth, keepaspectratio]{Chapter3_IMAGE/Activity_Diagram_quy_trình_xuất_bản_Manga.png}
    \caption{Activity Diagram quy trình xuất bản Manga}
    \label{fig:activity_xuat_ban}
\end{figure}
